\begin{frame}{왜 ``독립성''이 중요하냐}
  \begin{itemize}
    \item 미니배치 경사하강은 $\mathbb{E}[\nabla L]$을 \textbf{배치 평균}으로
      추정. 추정이 좋으려면(분산이 작으려면) 배치 안 샘플이
      \textbf{서로 독립}이어야 한다.
    \item 샘플이 거의 같으면 배치 평균은 \textbf{단 한 샘플의 경사}에 거의
      의존 $\Rightarrow$ 경사 추정치가 \textbf{극단적으로 높은 분산}.
  \end{itemize}
  \vskip0.5em
  \begin{block}{\kb{혼합시간}(mixing time)}
    ``32개를 모았다''는데 \textbf{효과적 표본수}는 32가 아니라 몇 개에
    불과. 연속 데이터가 ``몇 스텝을 건너뛰어야 무관해지는가''를
    재는 것이 혼합시간. CartPole 막대 각도는 몇십 스텝에 걸쳐 천천히
    바뀌므로 혼합시간도 \textbf{수십 스텝} $\Rightarrow$ 32개 연속 샘플은
    \textbf{동일 관찰 1개의 32번 반복}에 가깝다.
  \end{block}
\end{frame}
